\section{Related work}
\label{sec:related-work}

The closest comparisons concern transported complier effects, estimation of rough transport functionals, and inference under weak identification. Our immediate benchmark is \citet{CausalSmith2026TransportedLATEStrengthFrontier}, which characterizes asymptotically honest expected length through first-stage strength and transport-weight dispersion when the covariate and assignment geometry is supplied. Its feasible cell-based procedures learn target weights under uniform source cells and balanced assignment, or under regular nonuniform cells with supplied source probabilities and assignment propensities. Those results assume proportional source and target sample sizes and cell counts growing more slowly than the square root of the source sample size. The present comparison concerns the continuous-covariate model in \cref{obj:def:model-class}: the densities, encouragement propensity and arm means have Hölder exponent \leanref{sym:\beta}{\ensuremath{\beta}}, fixed at \(1/8\), and the two samples each contain \leanref{sym:n}{\ensuremath{n}} records.
% Source: [Honest Expected Length for Transported Complier Effects with Weak First Stages](https://causalsmith.org/papers/stat_transported_late_strength_frontier_v1/).

For identification and regular inference, \citet[Theorems 3.1, 4.2 and 4.3]{ChenHuang2025Generalizing} provide a particularly close comparison. Their observation scheme combines study records containing covariates, encouragement, receipt and outcomes with target covariates. Theorem 3.1 identifies the target complier average causal effect as a ratio of target-averaged source encouragement contrasts under their sampling overlap, target-population IV validity, and mean exchangeability restrictions for the outcome contrast and first stage. Their regular inference results distinguish supplied and estimated selection probabilities. Theorem 4.2 assumes known selection and assignment models, a fixed source-to-target sample-size ratio, and finite weighted second moments. Theorem 4.3 treats a correctly specified logistic selection model estimated by maximum likelihood, with known assignment probabilities and the stated moment and full-rank conditions. These theorems support delta-method inference at a fixed positive first stage. Our identification uses the transport restrictions on the complier-weighted outcome contrast and complier share in \cref{obj:ass:transport-complier-outcome,obj:ass:transport-complier-share}, while the inference criterion in \cref{obj:def:honest-intervals} requires coverage uniformly over the model class.
% Source: [Generalizing causal effects with noncompliance](https://arxiv.org/html/2506.00149).

A distinct observation scheme underlies \citet[Sections 5 and 7]{RudolphVanDerLaan2017Robust}. They transport outcomes to a site where covariates, encouragement and receipt are observed, so that the target first stage can be estimated directly. Their identification combines a common conditional outcome model across sites with target encouragement exchangeability, exclusion, monotonicity and positivity. Section 5 derives the transported-CACE influence curve and constructs a ratio of targeted estimators, with inference based on the influence curve or the multivariate delta method. Section 7 applies this construction to the Moving to Opportunity experiment. In the covariate-only target scheme studied here, both the outcome and receipt contrasts are obtained by averaging source contrasts over the target covariate distribution, as stated in \cref{obj:lem:transported-cace-identification}. This difference determines which transport restrictions identify the denominator as well as the numerator.
% Source: [Robust estimation of encouragement-design intervention effects transported across sites](https://pmc.ncbi.nlm.nih.gov/articles/PMC5784860/).

Transported principal-effect inference supplies another close comparison. Within their principal-stratification identification framework, \citet[Theorem 6]{Clark2024} establish asymptotic linearity, normality and efficiency for an influence-function estimator using trial assignment, receipt and outcome data together with target covariates. Their theorem assumes a positive limiting source-to-target sample-size ratio, nuisance estimation on a separate independent sample, uniform positivity of the relevant true and estimated probabilities, bounded outcome regressions, and consistency of the component influence functions. Efficient regular inference further requires each of the three specified products of nuisance estimation errors to be negligible relative to the inverse square root of sample size. These conditions describe a regular estimation regime for transported principal effects. The rough-class analysis here instead calibrates the observable estimator through the uniform mean-square bound in \cref{obj:lem:marked-cubic-mse}, which supplies the error control used by the score interval.
% Source: [Transportability of Principal Causal Effects](https://arxiv.org/html/2405.04419).

The reduced-form estimation scale has a direct methodological antecedent in \citet[Theorems 6 and 8]{ZengKennedyBodnarNaimi2025Efficient}. Their Theorem 6 gives a minimax lower bound for generalization when source and target covariates coincide, the joint covariate density for source membership and a specified treatment arm is constant, and the compound selection--treatment probability and outcome regression satisfy the stated positivity bounds. Its exponent specializes to one third for one-dimensional covariates and average nuisance smoothness one eighth. Their Theorem 8 gives conditional bias and variance bounds for a quadratic estimator of a transportation functional. That result assumes independent training data, uniform positivity, and uniformly well-conditioned true and estimated projection Gram matrices. Its bias bound includes both a product of projection approximation errors and a remainder multiplying two nuisance errors by the inverse-Gram estimation error. Thus, the exponent and the use of higher-order corrections are established antecedents, with the design, approximation and remainder conditions forming part of their results.
% Source: [Efficient Generalization and Transportation](https://arxiv.org/html/2302.00092).

The cubic construction in \cref{obj:def:cubic-estimator} belongs to the higher-order functional-estimation approach developed by \citet[Section 1]{Robins2008}. The marked densities in \cref{obj:def:marked-density-vector} express each transported contrast as the integral of a smooth rational function, as specified in \cref{obj:def:smooth-functional}. This representation permits a clipped pilot followed by observable linear, quadratic and cubic corrections computed from separate sample blocks. The construction specializes higher-order bias correction to the joint source--target observation scheme and provides the uniform mean-square control needed for finite-sample calibration.
% Source: [Higher order influence functions and minimax estimation of nonlinear functionals](https://arxiv.org/abs/0805.3040).

Score inversion connects that estimation problem to identification-robust inference. \citet[Theorem 3.2]{Ma2026IdentificationRobust} establish uniformly correct asymptotic size for an orthogonalized test of a LATE moment under their high-dimensional nuisance, sparsity, overlap, moment and empirical regularity conditions. The result accommodates weak identification within that class. Our score interval in \cref{obj:def:score-interval} uses a finite-sample calibration derived from a uniform mean-square bound over the specified rough class. The shared principle is to test an affine moment in the candidate effect; the calibration and the model over which coverage is required determine the resulting guarantee.
% Source: [Identification-robust inference for the LATE with high-dimensional covariates](https://arxiv.org/html/2302.09756).

The geometry of this inversion is clarified by \citet[Section 3 and Theorems 1 and 2]{SmuclerLanniMasip2025Score}. Their studentized score inversion reduces to a quadratic inequality and can produce bounded or unbounded intervals, disconnected sets, the whole line, an empty set or a singleton. Under their independent nuisance estimation, boundedness, overlap, consistency and product-rate conditions, Theorem 1 shows agreement with the doubly robust Wald interval up to stochastic endpoint errors of order \(O_P(n^{-1})\), with probability tending to one, at a fixed law with a nonzero first stage. Theorem 2 gives a nonnormal ratio limit under their local-to-zero moment and nuisance-negligibility conditions. In our construction, a deterministic calibration radius and the bounded causal domain yield the connected acceptance interval and fallback rules specified in \cref{obj:def:score-interval}. These comparisons locate the inversion principle within established weak-identification methodology.
% Source: [A note on the properties of the confidence set for the local average treatment effect obtained by inverting the score test](https://arxiv.org/html/2506.10449).

Two older principles organize the criterion used here. Anderson--Rubin inversion constructs confidence sets by testing an affine reduced-form moment without dividing by a possibly weak first stage \citep[pp. 46--63]{Anderson1949}; Fieller's ratio analysis likewise allows denominator uncertainty to determine confidence-set geometry \citep[pp. 175--185]{Fieller1954}. The present interval specializes that logic to the bounded causal domain and replaces an asymptotic critical value with the finite-sample rough-functional radius in \cref{obj:def:score-interval}. On the decision-theoretic side, weak identification can force unbounded confidence sets on an unrestricted parameter space \citep[pp. 1365--1387]{Dufour1997}, while honest expected length over bounded functional classes is governed by between-class modulus arguments \citep[pp. 1805--1840]{CaiLow2004Adaptation}. Here bounded outcomes cap length at constant order, and the legal transported-IV family supplies the testing comparison that yields the second branch of the frontier.

The contribution is the joint equal-size, rough-geometry, uniform-coverage and expected-length characterization in \cref{obj:thm:sharp-length-frontier-full-elbow}. Its upper bound uses an observable cubic estimator and a calibrated score interval; its lower bound is realized by the primitive IV family in \cref{obj:def:legal-iv-mixture,obj:lem:legal-iv-mixture-full-elbow}. That family retains the causal and transport restrictions of the model while producing the observed experiments needed for the length comparison. Together, these results connect the established rough-functional exponent and score-inversion principle to a sharp expected-length criterion for transported complier effects with unknown continuous covariate and assignment geometry.
